\documentclass[a4paper,11pt]{article}
\usepackage[utf8]{inputenc}
\usepackage[T1]{fontenc}
\usepackage[french]{babel}
\usepackage[left=2cm,right=2cm,top=2cm,bottom=2cm]{geometry}
\usepackage{xcolor}
\usepackage{hyperref}
\usepackage{fontawesome5}
\usepackage{parskip}

% --- Couleurs CEA ---
\definecolor{primary}{RGB}{0, 102, 153}
\hypersetup{colorlinks=true, urlcolor=primary}

\begin{document}

% --- EXPÉDITEUR ---
\noindent
\begin{minipage}[t]{0.5\textwidth}
    \textbf{Loïc Aron MBASSI EWOLO}\\
    Élève-Ingénieur Bac+5 -- Double diplôme ENSEA\\[3pt]
    \faMapMarker\ Cergy, France\\
    \faPhone\ (+33) 7 59 52 33 98\\
    \faEnvelope\ \href{mailto:loicaron.mbassiewolo@ensea.fr}{loicaron.mbassiewolo@ensea.fr}
\end{minipage}
\hfill
% --- DATE ---
\begin{minipage}[t]{0.4\textwidth}
    \raggedleft
    Cergy, le 27 février 2026
\end{minipage}

\vspace{1.2cm}

% --- DESTINATAIRE ---
\hfill
\begin{minipage}[t]{0.5\textwidth}
    \raggedleft
    \textbf{CEA LITEN -- Institut Liten}\\
    Service Recrutement\\
    Grenoble, France
\end{minipage}

\vspace{1cm}

\noindent\textbf{Objet :} Candidature -- Alternance Développement Logiciel ACV (Python / MyLCA) -- Dès Septembre 2026

\vspace{0.8cm}

Madame, Monsieur,

Le développement de MyLCA -- un outil Python modulaire, paramétrique et open source pour l'Analyse de Cycle de Vie -- est précisément le type de projet qui se situe à l'intersection de ma formation en génie logiciel et de mes expériences en développement d'outils scientifiques. C'est cette adéquation, doublée de l'ancrage du CEA LITEN dans la recherche sur les technologies de l'énergie, qui motive ma candidature à l'alternance proposée.

Un système multi-agents en Python que j'ai conçu pour des recommandations agricoles (FastAPI, Google ADK, RAG, Gemini) m'a confronté directement aux enjeux centraux de votre offre : structurer des bibliothèques de modèles réutilisables, orchestrer des sources de données hétérogènes, et exposer des fonctionnalités complexes à des utilisateurs non-développeurs via des interfaces simplifiées. L'architecture modulaire et le déploiement Docker pour la reproductibilité de ce projet -- au même titre que l'approche open source de Brightway2 -- ont structuré ma manière d'aborder l'ingénierie logicielle orientée recherche. Par ailleurs, mon outil de génération de code Laravel (open source, +30 étoiles GitHub, déployé sur Packagist) témoigne de ma capacité à construire, maintenir et itérer sur un logiciel complexe utilisé par d'autres.

Ma formation en mathématiques appliquées (algèbre linéaire, probabilités, analyse numérique) et mon expérience en tant qu'assistant de recherche au MILA (Institut Québécois d'Intelligence Artificielle, été 2025) me donnent les bases pour aborder les aspects quantitatifs de l'ACV prospective et dynamique : travail sur des simulations Python itératives, analyse de convergence, reproduction d'expériences avec rigueur. Cette posture de recherche -- où la qualité du code conditionne la robustesse des résultats -- me semble directement applicable aux études de sensibilité et aux analyses Monte-Carlo que décrit l'offre. Enfin, un projet de traitement automatisé d'électrocardiogrammes (Python, TensorFlow, collaboration avec mentors Google DeepMind) a renforcé ma pratique du calcul scientifique appliqué à des données complexes.

Je suis disponible dès septembre 2026 pour rejoindre les équipes ACV du LITEN à Grenoble. Contribuer au déploiement de MyLCA dans des laboratoires de recherche sur l'énergie représente une opportunité concrète de mettre mes compétences logicielles au service d'un enjeu environnemental structurant.

Je reste à votre disposition pour un entretien afin de discuter de ma candidature.

Je vous prie d'agréer, Madame, Monsieur, l'expression de mes salutations distinguées.

\vspace{0.8cm}

\textbf{Loïc Aron MBASSI EWOLO}\\
\textit{Élève-Ingénieur ENSEA / Polytechnique Yaoundé}

\end{document}
